\documentclass[11pt]{article}
\usepackage[margin=1in]{geometry}
\usepackage{textcomp}
\usepackage{amsmath,amssymb,amsthm,booktabs,hyperref}
\hypersetup{colorlinks=true,linkcolor=black,citecolor=black,urlcolor=blue}
\newtheorem{theorem}{Theorem}
\newtheorem{lemma}[theorem]{Lemma}
\newtheorem{proposition}[theorem]{Proposition}
\theoremstyle{remark}
\newtheorem{remark}[theorem]{Remark}
\newcommand{\R}{\mathbb{R}}
\title{Certified lower bounds for $\lambda_1\cdot\mathrm{Area}$ on cells of pentagon shape space\\
by piecewise-affine transplantation\\[4pt]
\large Computer-assisted proof candidate (for the stated cells), not peer reviewed}
\author{Puyi Sun\thanks{ORCID 0009-0004-1644-5300. See Section~\ref{sec:ai} for the roles of AI systems in this work.}}
\date{__DATE__ \quad (version __VERSION__)}
\begin{document}
\maketitle

\begin{abstract}
For a planar domain $\Omega$ let $F(\Omega)=\lambda_1(\Omega)\,|\Omega|$, where $\lambda_1$ is the first Dirichlet
eigenvalue of the Laplacian. The P\'olya--Szeg\H{o} conjecture asserts that among pentagons the regular one minimises $F$;
its value is $F_*\approx 18.9191043$. We describe a method that certifies $F(P)>F_*$ simultaneously for \emph{every}
pentagon $P$ in a box (``cell'') of a six-dimensional shape chart, using only the first eigenvalues of the $64$ corner
pentagons of the box and explicit matrix inequalities obtained from a piecewise-affine transplantation onto a fixed
reference fan. Everything is evaluated in ball arithmetic. We present six certificates for five distinct parameter boxes
at three centre locations (two certificates apply different methods to the same box), including the cell identified as
hardest in a publicly available outward-continuation prototype \cite{liu}, and a partial independent second-party replay
of certificate components (Section~\ref{sec:replay}): __REPLAY_ABSTRACT__ This is a method and a set of cell certificates; it is \emph{not} a proof of
the conjecture, nor of any statement about a neighbourhood of the regular pentagon.
\end{abstract}

\section{Background}
P\'olya and Szeg\H{o} proved that the equilateral triangle and the square minimise $F$ among triangles and
quadrilaterals, $F\ge 4\pi^2/\sqrt3$ and $F\ge 2\pi^2$ respectively, and conjectured the same for $n$-gons with
$n\ge5$ \cite{ps}. Bogosel and Bucur reduced the case of fixed $n$ to a finite (if large) computation and proved local
minimality of the regular pentagon and hexagon with validated computing \cite{bb-jep,bb-val}; local minimality of the
regular heptagon and octagon was certified in \cite{ep}; an explicit local ball around the regular pentagon and a
prototype of an outward continuation were published in \cite{liu}; the conjecture is known for convex polygons with
sufficiently many sides \cite{cghl}. To our knowledge (two independent literature searches, October 2026) the global
pentagon case remains open. The present note contributes a cell-certification tool for the region away from the
regular pentagon; we make no claim of priority beyond the specific method and certificates described here.

\section{Chart, reference fan and transplantation}
A pentagon is $P(x)$ with vertices $V_0=(0,0)$, $V_1=(1,0)$, $V_2=(x_1,x_2)$, $V_3=(x_3,x_4)$, $V_4=(x_5,x_6)$,
counter-clockwise. $F$ is invariant under similarities. Let $P_0$ be a fixed regular pentagon (binary64 vertices taken
exactly) with vertex centroid $c_0$, and $T_k=\mathrm{conv}(c_0,P_{0,k},P_{0,k+1})$, $E_k=[P_{0,k}-c_0,\;P_{0,k+1}-c_0]$.
For $P=P(x)$ with vertex centroid $c$ put
\[
D_k=[V_k-c,\;V_{k+1}-c]\,E_k^{-1},\qquad J_k=\det D_k,\qquad G_k=J_k\,D_k^{-1}D_k^{-T}\qquad(k=0,\dots,4).
\]
If $P$ is strictly convex, $c$ is interior, every $J_k>0$ and the piecewise-affine map $\Phi$ equal to
$y\mapsto c+D_k(y-c_0)$ on $T_k$ is a bi-Lipschitz homeomorphism $P_0\to P$; then
\begin{equation}\label{eq:rq}
\lambda_1(P)=\inf_{0\ne v\in H^1_0(P_0)}\frac{\sum_k\int_{T_k}\nabla v^{T}G_k\nabla v}{\sum_k J_k\int_{T_k}v^2},
\qquad |P|=\sum_k |T_k|\,J_k .
\end{equation}
\emph{Convexity gate.} The $15$ determinants $\det(V_{i+1}-V_i,\,V_j-V_i)$, $j\notin\{i,i+1\}$, are affine in each
coordinate of $x$; positivity at the $64$ corners of a box therefore implies strict convexity on the whole box.
(Positivity of the $J_k$ alone does not suffice: the regular pentagram is a counterexample.)

\section{The cell inequality}
A cell is a box $C=\prod_i[\ell_i,h_i]$ with binary64 endpoints. Write $x(u)=\ell+(h-\ell)\odot(1+u)/2$,
$u\in[-1,1]^6$, let $x_c$ ($c\in\{\pm1\}^6$) be the corners and $\alpha_c(u)=\prod_i(1+c_iu_i)/2$ the product
(multilinear) weights; $D_{c,k},J_{c,k},G_{c,k}$ denote the corner quantities.

\begin{lemma}[exact interpolation]\label{lem:interp}
$D_k(x(u))=\sum_c\alpha_c(u)D_{c,k}$ and $J_k(x(u))=\sum_c\alpha_c(u)J_{c,k}$.
\end{lemma}
\begin{proof}
$D_k$ is affine in $x$ and multilinear interpolation reproduces affine functions. $J_k$ is a quadratic polynomial in
$x$ whose pure squares vanish: the derivative of $D_k$ in a single coordinate $x_i$ moves one coordinate of one vertex
and has rank $\le1$, so its determinant is $0$; hence $J_k$ is multiaffine and equals its multilinear interpolant.
\end{proof}

\begin{lemma}[covariance identity]\label{lem:cov}
With $X_c=D_{c,k}^{-1}$, $w_c=\alpha_cJ_{c,k}/J_k$, $M=\sum_cw_cX_c$ and $S_k=\sum_c\alpha_cG_{c,k}$,
\[
G_k=S_k-J_k\,\mathcal V_k,\qquad \mathcal V_k=\sum_cw_c(X_c-M)(X_c-M)^T\succeq0 .
\]
\end{lemma}
\begin{proof}
For $2\times2$ matrices the adjugate is linear, so $\mathrm{adj}\,D_k=\sum_c\alpha_c\,\mathrm{adj}\,D_{c,k}=J_k M$;
then $G_k=\mathrm{adj}D_k(\mathrm{adj}D_k)^T/J_k=J_kMM^T$ and $\sum_c\alpha_cG_{c,k}=J_k\sum_cw_cX_cX_c^T$.
\end{proof}

\begin{lemma}[matrix Efron--Stein]\label{lem:es}
Let $G_l\succ0$ satisfy $G_l\preceq G_{c,k}$ for all corners, $J^{\max}_k=\max_cJ_{c,k}$, $p_iq_i=(1-u_i^2)/4$, and for
each coordinate $i$ let $\Delta_i^{(e)}$ be the $32$ edge differences $X_c-X_{c'}$ ($c,c'$ differing in coordinate $i$),
$R_i$ their mean and $d^{(e)}_i=\Delta^{(e)}_i-R_i$. Then for every $\tau>0$
\[
J_k\mathcal V_k\preceq \epsilon_k(u)\,G_l\preceq\epsilon_k(u)\,S_k,
\]
\begin{multline*}
\epsilon_k(u)=J^{\max}_k\Big[(1+\tau)\,\lambda_{\max}\Big(G_l^{-1/2}\sum_ip_iq_iR_iR_i^TG_l^{-1/2}\Big)\\
+\big(1+\tfrac1\tau\big)\sum_ip_iq_i\max_e\lambda_{\max}\big(d_i^{(e)T}G_l^{-1}d_i^{(e)}\big)\Big].
\end{multline*}
\end{lemma}
\begin{proof}
$J_k\mathcal V_k\preceq\sum_c\alpha_cJ_{c,k}(X_c-\bar X)(X_c-\bar X)^T\preceq J^{\max}_k\,\mathrm{Var}_\alpha(X)$ for the
$\alpha$-mean $\bar X$. The Efron--Stein inequality for the product measure $\alpha$, applied to every scalar projection,
gives the Loewner bound $\mathrm{Var}_\alpha(X)\preceq\sum_ip_iq_i\,\mathbb E_\alpha[\Delta_i\Delta_i^T]$, and
$(R+d)(R+d)^T\preceq(1+\tau)RR^T+(1+\tau^{-1})dd^T$ \cite{es}.
\end{proof}

\emph{Choice of $G_l$.} For each sector let $G_{\rm cen}=\frac1{64}\sum_cG_{c,k}$. Two admissible candidates are computed
in ball arithmetic: (a) $G_{\rm cen}-\eta I$ with $\eta\ge\max_c\|G_{c,k}-G_{\rm cen}\|_2$, if it is positive definite;
(b) $\mu_lG_{\rm cen}$ with $0<\mu_l\le\min_c\lambda_{\min}(G_{\rm cen}^{-1/2}G_{c,k}G_{\rm cen}^{-1/2})$ (the smallest root
of $\det(G_{c,k}-\mu G_{\rm cen})=0$, rounded down). Both satisfy $G_l\preceq G_{c,k}$ for every corner; the certificate
uses the one with the larger $\lambda_{\min}(G_{\rm cen}^{-1}G_l)$ (a choice, not a bound), and asserts $G_l\succ0$.
($\mu_l$ is unrelated to the quantity $\mu$ of Section~4.) The value of $\tau$ is the best of a fixed finite list.

\begin{theorem}[cell bound]\label{thm:cell}
Let $F^L_c\le F(P(x_c))$ be certified corner lower bounds, $\lambda^L_c=F^L_c/|P(x_c)|$, and suppose the convexity gate
holds on $C$ and $\epsilon_k(u)<1$. Then for every $u\in[-1,1]^6$
\begin{equation}\label{eq:t1}
F(P(x(u)))\;\ge\;\Big(1-\max_k\epsilon_k(u)\Big)\,\min_k\,|P(x(u))|\,\frac{\sum_c\alpha_c(u)\,\lambda^L_c\,J_{c,k}}{J_k(x(u))} .
\end{equation}
\end{theorem}
\begin{proof}
By Lemmas~\ref{lem:cov}--\ref{lem:es}, $G_k\succeq(1-\epsilon_k)S_k$. For $v\in H^1_0(P_0)$ the numerator of
\eqref{eq:rq} is $\ge\sum_k(1-\epsilon_k)\sum_c\alpha_cN_{c,k}(v)\ge(1-\max_k\epsilon_k)\sum_c\alpha_cN_c(v)$, where
$N_c=\sum_kN_{c,k}$ is the numerator for the corner $P(x_c)$, and $N_c(v)\ge\lambda_c\sum_kJ_{c,k}\int_{T_k}v^2$.
The denominator is $\sum_kJ_k\int_{T_k}v^2$ by Lemma~\ref{lem:interp}; a ratio of positive sums is at least the smallest
termwise ratio. (The corner Rayleigh inequality holds only for the sum over sectors, which is why the factor is
$1-\max_k\epsilon_k$ and not paired sector by sector.)
\end{proof}

\emph{Exact-covariance alternative.} By Lemma~\ref{lem:cov}, $S_k-G_k\succeq0$, so
\[
e_k(u)=\lambda_{\max}\big(S_k^{-1/2}(S_k-G_k)S_k^{-1/2}\big)\in[0,1)
\]
satisfies $G_k\succeq(1-e_k)S_k$ directly; here $S_k$ and $\mathrm{adj}\,D_k$ (hence
$G_k=\mathrm{adj}D_k(\mathrm{adj}D_k)^T/J_k$) are multiaffine. Any upper bound of
$\epsilon_k$ or $e_k$ may be used in place of $\epsilon_k$ in Theorem~\ref{thm:cell}.

Two ways of using \eqref{eq:t1} are implemented.

\emph{Min-corner form.} Write $f_k=J_k/|P|$ and $I(f_k)=\sum_c\alpha_c f_k(x_c)$, so that with
$\lambda^L_c\ge F_{\min}/|P(x_c)|$, $F_{\min}=\min_cF^L_c$, the $k$-th term of \eqref{eq:t1} is at least
$F_{\min}\,I(f_k)/f_k$. Tensor-linear interpolation gives $|I(f_k)-f_k|\le\mathrm{err}_k=\sum_i\frac{(h_i-\ell_i)^2}{8}
\sup_C|\partial_i^2f_k|$ with the full side lengths $h_i-\ell_i$, the supremum being enclosed on $64$ sub-boxes whose
union is $C$ (interval evaluation, asserted coverage). With $f_k^{\min}$ a positive lower bound of $f_k$ on $C$ and
$\epsilon_k$ evaluated at $p_iq_i=1/4$ (its maximum over $u$),
\[
F(P)\ \ge\ F_{\min}\cdot\min_k(1-\epsilon_k)\cdot\Big(1-\max_k\frac{\mathrm{err}_k}{f_k^{\min}}\Big)=:\mathrm{factor}\cdot F_{\min}
\quad\text{on all of }C,
\]
provided all $J_{c,k}>0$, $f_k^{\min}>0$, $\epsilon_k<1$ and $\mathrm{err}_k<f^{\min}_k$ (all asserted). The code also
multiplies by $J_{\min}/(J_{\min}+\delta)$ with $\delta=\sum_i(\frac{h_i-\ell_i}2)^2|\det\partial_iD_k|$; by
Lemma~\ref{lem:interp} $\delta=0$, so this factor is $1$.

\emph{T1 form.} Since $Q_k>0$ and $\epsilon<1$, the right-hand side of \eqref{eq:t1} equals
$\min_{k,k'}Q_k(u)\,(1-\epsilon_{k'}(u))$ with $Q_k=|P|\sum_c\alpha_c\lambda^L_cJ_{c,k}/J_k$. It is bounded below over
$u\in[-1,1]^6$ by branch and bound: on a sub-box with centre $c$ and radii $r$, each of the $25$ products is bounded below
by $f(c)-\sum_i|\partial_if(c)|r_i-\frac12\sum_{i,j}\sup_{\rm box}|\partial_{ij}f|\,r_ir_j$ (point jet at the centre,
interval Hessian over the box); multiaffine pieces have exact ranges from the sub-box corner values. For $\epsilon_k$ and
$e_k$ the $\lambda_{\max}$ of a positive semidefinite $2\times2$ matrix is replaced by the polynomial majorant
$\mathrm{tr}/2+(q+s^2)/(2s)$, $q=\mathrm{tr}^2/4-\det$ (AM--GM, $s>0$ fixed per box), or by the trace. Cells C, D and F
use the Efron--Stein bound only; cell E uses, per sub-box and sector, the smaller of the valid majorants of $\epsilon_k$
and $e_k$. The accepted leaves form an exact dyadic partition of $[-1,1]^6$ and are part of the certificate.

\section{Corner eigenvalues and the regular pentagon}
Corner lower bounds use the method of particular solutions: $u=\sum_{j,m}c_{jm}J_{\nu_{jm}}(kr_j)\sin(\nu_{jm}\theta_j)$,
$\nu_{jm}=m\pi/\alpha_j$, solves $\Delta u+\tilde\lambda u=0$ exactly ($k^2=\tilde\lambda$); by Fox--Henrici--Moler some
eigenvalue satisfies $|\tilde\lambda-\lambda_n|\le\mu\lambda_n$ with $\mu=|P|^{1/2}\sup_{\partial P}|u|/\|u\|_2$
\cite{fhm,mp}; $\|u\|_2$ is bounded below through the Helmholtz mean-value identity on an inscribed disk; and, when
$\mu<1$, $\tilde\lambda|P|/(1-\mu)<2\pi j_{0,1}^2\le\lambda_2|P|$ (Hong--Krahn--Szeg\H{o} \cite{krahn,hong,henrot}) identifies $\lambda_n=\lambda_1$.
Hence $\lambda_1\ge\tilde\lambda/(1+\mu)$. All $384$ corners have $\mu\le3.1\cdot10^{-4}$. Every corner certificate
records the vertices, $\tilde\lambda$ (a binary64 number), $\mu$ and the $40$ coefficients, so that a verifier evaluates
the \emph{same} explicit function. The producer verifier treats the stored enclosures of $\mu$ and of the area as
certificate data: it re-checks $\mu<1$, the index inequality and $\tilde\lambda/(1+\mu)$ from them, but does not
recompute the boundary supremum or the disk bound; those are recomputed by the second-party checker for the corners
listed in Section~\ref{sec:replay}. All corner pentagons are strictly convex (convexity gate), so every $\theta_j$ is
taken in $[0,\alpha_j]$ with $\alpha_j<\pi$ and the branch cut of $\theta_j$ lies outside $\overline P$. The stored
binary64 vertex coordinates, $\tilde\lambda$ and coefficients define the trial function exactly; interval arithmetic
encloses its evaluations, and the final reported lower and upper bounds are rounded outward to binary64.

The regular pentagon is bounded from above in two independent ways:
(i) the FHM upper bound $\tilde\lambda/(1-\mu)$ at a binary64 regular pentagon $P_t$ whose copy shrunk by the factor
$1-10^{-12}$ lies inside the exact regular pentagon (checked in ball arithmetic), with domain monotonicity, giving
$F_*\le 18.919134195733918$; (ii) a conforming $P_2$ finite-element Rayleigh quotient on an inscribed rational
pentagon with directed-rounding exact integration (second party; source and logs in \texttt{second\_party/p2\_fem}),
giving $F_*\le 18.919118220170337$. The trial vector of that logged run is not archived, so this value is a
historical log that cannot be replayed digit for digit from the package; rerunning the two scripts certifies a fresh
upper bound whose last digits may differ. All margins in this note are computed against the larger value (i),
$F_*^{\rm up}=18.919134195733918$. Validated eigenvalue bounds on polygons by other means are available in \cite{lo}.

\section{Results}\label{sec:results}
__RESULTS_TABLE__

\section{Cost of a cover (estimate)}\label{sec:cost}
__COST_SECTION__

\section{Independent second-party replay}\label{sec:replay}
__REPLAY_TABLE__

\section{Limitations}
The certified boxes are small (side $2\rho\le0.06$) at chart distance $0.1$--$0.57$ from the regular pentagon; no statement is made about the region between them,
about a neighbourhood of the regular pentagon, about non-convex pentagons, or about the conjecture. Near the regular
pentagon the gap $F-F_*$ shrinks quadratically and the corner-eigenvalue accuracy, not the transplantation bound,
limits the cell size. The arithmetic relies on Arb/python-flint \cite{arb} (producer) and on mpmath interval arithmetic
and Python's \texttt{Fraction}/\texttt{Decimal} (second party); no formal proof assistant was used.

\section{Roles of AI systems}\label{sec:ai}
Claude (Anthropic) designed the method, wrote and revised the producer algorithms and code, and interpreted the
resulting logs. GPT (OpenAI) did the literature review and the mathematical audit and wrote the independent
second-party checkers (Section~\ref{sec:replay}), whose code shares no source with the producer. The code ran as
local Python processes, in the AI systems' sandboxes and on the author's computer; the author initiated and
coordinated the runs and decided on the release. The author did not establish the mathematical correctness by hand;
responsibility for the publication and its correctness lies with the author. Errors, if any, are reported in the
repository and corrected in later versions.

\section*{Data and code}
All certificates, producer code, independent checkers, replay commands and a SHA-256 manifest are in the accompanying
repository (CC BY 4.0 for text and data, MIT for code).

\begin{thebibliography}{99}
\bibitem{ps} G.~P\'olya, G.~Szeg\H{o}, \emph{Isoperimetric Inequalities in Mathematical Physics}, Princeton, 1951.
\bibitem{bb-jep} B.~Bogosel, D.~Bucur, On the polygonal Faber--Krahn inequality, \emph{J. \'Ec. polytech. Math.} 11 (2024) 19--105, doi:10.5802/jep.250.
\bibitem{bb-val} B.~Bogosel, D.~Bucur, Polygonal Faber--Krahn inequality: local minimality via validated computing, arXiv:2406.11575.
\bibitem{ep} Evidence Press, Certified local minimality of the regular heptagon and octagon, release v1.0.1, doi:10.5281/zenodo.23066520 (2026), not peer reviewed.
\bibitem{liu} X.~Liu, \texttt{github.com/xfliu/polya-szego-conjecture} (public release, September 2026).
\bibitem{cghl} Z.~Cheng, C.~Gui, Y.~Hu, Q.~Li, The P\'olya--Szeg\H{o} conjecture for convex polygons with many sides, arXiv:2609.18500 (2026).
\bibitem{fhm} L.~Fox, P.~Henrici, C.~Moler, Approximations and bounds for eigenvalues of elliptic operators, \emph{SIAM J. Numer. Anal.} 4 (1967) 89--102, doi:10.1137/0704008.
\bibitem{mp} C.~B.~Moler, L.~E.~Payne, Bounds for eigenvalues and eigenvectors of symmetric operators, \emph{SIAM J. Numer. Anal.} 5 (1968) 64--70, doi:10.1137/0705004.
\bibitem{krahn} E.~Krahn, \"Uber Minimaleigenschaften der Kugel in drei und mehr Dimensionen, \emph{Acta Comm. Univ. Tartu (Dorpat)} A9 (1926) 1--44.
\bibitem{hong} I.~Hong, On an inequality concerning the eigenvalue problem of membrane, \emph{K\={o}dai Math. Sem. Rep.} 6 (1954) 113--114.
\bibitem{henrot} A.~Henrot, \emph{Extremum Problems for Eigenvalues of Elliptic Operators}, Birkh\"auser, 2006, Chapter~4.
\bibitem{lo} X.~Liu, S.~Oishi, Verified eigenvalue evaluation for the Laplacian over polygonal domains of arbitrary shape, \emph{SIAM J. Numer. Anal.} 51 (2013) 1634--1654, doi:10.1137/120878446.
\bibitem{es} B.~Efron, C.~Stein, The jackknife estimate of variance, \emph{Ann. Statist.} 9 (1981) 586--596.
\bibitem{arb} F.~Johansson, Arb: efficient arbitrary-precision midpoint-radius interval arithmetic, \emph{IEEE Trans. Comput.} 66 (2017) 1281--1292, doi:10.1109/TC.2017.2690633.
\end{thebibliography}
\end{document}
